%% ----------------------------------------------------------------------------
%%  §11.3  Serie a termini non negativi
%% ----------------------------------------------------------------------------
\section{Serie a termini non negativi}
\label{sec:serie-positive}

Diremo che una serie $\sum a_{k}$ è \emph{a termini non negativi} se $a_{n}\ge0$ per ogni $n$,
e \emph{a termini positivi} se $a_{n}>0$ per ogni $n$. Per queste serie la situazione è
particolarmente semplice, perché la successione delle somme parziali è monotòna: da
$a_{n+1}\ge0$ segue
\begin{equation}\label{eq:somme-crescenti}
  s_{n+1}=s_{n}+a_{n+1}\ge s_{n} \qquad\perogni n\in\N .
\end{equation}

\begin{teorema}[sulle serie a termini non negativi]\label{teo:serie-positive}
\index{serie!a termini non negativi}
Una serie a termini non negativi non può essere indeterminata: è convergente oppure
divergente positivamente. Più precisamente, è convergente se e solo se la successione
delle sue somme parziali è limitata, e in tal caso
\begin{equation}\label{eq:somma-sup}
  \sum_{k=1}^{\infty}a_{k}=\sup_{n\in\N}s_{n} .
\end{equation}
\end{teorema}

\needspace{5\baselineskip}
\begin{schema}
\Ipotesi{$a_{n}\ge0$ per ogni $n$.}
\Tesi{$s_{n}$ ha limite, finito o $+\infty$; \ $\sum a_{k}<+\infty$ $\sse$ $s_{n}$ limitata; \
$\sum a_{k}=\sup s_{n}$.}
\end{schema}

\begin{strategia}
Per la \eqref{eq:somme-crescenti} la successione $s_{n}$ è crescente, e il teorema sulle
successioni monotòne dice tutto: il limite esiste ed è l'estremo superiore, finito se e
solo se la successione è limitata superiormente. Inferiormente $s_{n}$ è limitata da $s_{1}$.
\end{strategia}

\begin{dimostrazione}
Per la \eqref{eq:somme-crescenti} la successione $s_{n}$ è crescente; per il teorema sulle
successioni monotòne (\cref{teo:monotone}) essa ha limite, e $\lim s_{n}=\sup_{n}s_{n}$.
Il limite non può essere $-\infty$, perché $s_{n}\ge s_{1}$ per ogni $n$, e quindi la serie è
convergente o divergente a $+\infty$. Inoltre $\sup_{n}s_{n}<+\infty$ se e solo se $s_{n}$ è
limitata superiormente, cioè (essendo $s_{n}\ge s_{1}$) limitata.
\end{dimostrazione}

In particolare, \emph{la somma di una serie convergente a termini non negativi è
maggiore o uguale a ciascuna delle sue somme parziali}. Per una serie a termini non
negativi si scrive spesso $\sum a_{k}<+\infty$ per dire che la serie converge, e la
condizione necessaria (\cref{teo:condizione-necessaria}) diventa un criterio di
divergenza: se $a_{n}$ non tende a zero, la serie diverge a $+\infty$.

\begin{esempio}\label{es:k-su-k1}
La serie $\sum_{k=1}^{\infty}\frac k{k+1}$ è a termini positivi, quindi convergente o divergente
a $+\infty$. Non è convergente, perché il suo termine generale tende a $1$ e non a $0$.
Dunque $\sum_{k=1}^{\infty}\frac k{k+1}=+\infty$.
\end{esempio}

\begin{esempio}[La serie dei reciproci dei fattoriali]\label{es:serie-fattoriali}
Consideriamo la serie a termini positivi $\sum_{k=0}^{\infty}\frac1{k!}$ (con $0!=1$). Per $k\ge1$
si ha $k!=1\cdot2\cdots k\ge2^{k-1}$, perché ciascuno dei $k-1$ fattori $2,3,\dots,k$ è almeno
$2$. Quindi, per la formula della progressione geometrica (\cref{prop:geometrica}),
\[
  s_{n}=1+\sum_{k=1}^{n}\frac1{k!}\le1+\sum_{k=1}^{n}\frac1{2^{k-1}}=1+\frac{1-2^{-n}}{1-\frac12}<3 :
\]
le somme parziali sono limitate, e la serie converge, con somma compresa tra $s_{1}=2$ e
$3$. La stima del resto \eqref{eq:stima-e} della formula di Taylor dice di più:
$0<e-\sum_{k=0}^{n}\frac1{k!}<\frac3{(n+1)!}\to0$, e quindi
\begin{equation}\label{eq:e-serie}
  \sum_{k=0}^{\infty}\frac1{k!}=1+1+\frac12+\frac1{3!}+\frac1{4!}+\dots=e .
\end{equation}
Il teorema sulle serie a termini non negativi ha garantito l'\emph{esistenza} della
somma con un ragionamento elementare; la formula di Taylor ne ha fornito il
\emph{valore}. (Lo stesso valore si ottiene anche senza Taylor, direttamente dalla
definizione di $e$ come limite di $\bigl(1+\frac1n\bigr)^{n}$: si veda
l'\cref{oss:e-due-vie}.) È il metodo efficiente per il calcolo di $e$ annunciato nel
\cref{cap:successioni}: undici termini danno sette cifre decimali esatte, mentre la
successione $\bigl(1+\frac1n\bigr)^{n}$ ne dà appena cinque con $n=10^{6}$.
\end{esempio}
